\documentclass[11pt]{article}
\usepackage[margin=1in]{geometry}
\usepackage{amsmath}
\usepackage{amssymb}
\usepackage{hyperref}

\title{Coordinate System \& Symmetry Specification\\for Warp-Bubble Metrics}
\author{Arcticoder}
\date{May 29, 2025}

\begin{document}
\maketitle

\section{Coordinate Chart}
We adopt a standard spherical coordinate system on the background spacetime:
\begin{align*}
  &\mathrm{Coordinates:}\quad (t,\,r,\,\theta,\,\phi) \\
  &\mathrm{Domains:}\quad
    t \in (-\infty, +\infty),\quad
    r \in [0, +\infty),\quad
    \theta \in [0,\pi],\quad
    \phi \in [0,2\pi).
\end{align*}
Here $t$ is the time coordinate, $r$ the radial coordinate, and $(\theta,\phi)$ the usual polar and azimuthal angles.

\section{Symmetry Assumptions}
To simplify the metric ansatz for warp-bubble geometries, we impose:
\begin{enumerate}
  \item \textbf{Axial symmetry} about the $z$-axis: \\
    All metric components are independent of the azimuthal angle $\phi$.
  \item \textbf{Equatorial reflection symmetry}: \\
    Invariance under $\theta \mapsto \pi - \theta$.
  \item \textbf{Compact radial support}: \\
    Nontrivial warp functions only for $r \le R$, where $R$ is the bubble radius.  Outside this region the spacetime matches the background (e.g.\ Minkowski).
\end{enumerate}

\section{General Static, Axisymmetric Metric}
A fully general static, axisymmetric line element in these coordinates would include off–diagonal terms and $\theta$–dependence:
\[
  ds^2 =
    -\,g_{tt}(r,\theta)\,dt^2
    +2\,g_{t\phi}(r,\theta)\,dt\,d\phi
    +g_{rr}(r,\theta)\,dr^2
    +2\,g_{r\theta}(r,\theta)\,dr\,d\theta
    +g_{\theta\theta}(r,\theta)\,d\theta^2
    +g_{\phi\phi}(r,\theta)\,d\phi^2.
\]

\section{Reduced Metric Ansatz}
Imposing staticity, axial and reflection symmetry, and compact support in $r$, all off–diagonal pieces vanish and the remaining metric functions depend only on $r$.  We therefore arrive at:
\[
  \boxed{
    ds^2 = -A(r)\,dt^2
      + B(r)\,dr^2
      + C(r)\,r^2\,d\theta^2
      + D(r)\,r^2\sin^2\theta\,d\phi^2
  }
\]
where
\[
  A(r),\,B(r),\,C(r),\,D(r)
  \;\begin{cases}
    \equiv 1, & r > R \quad (\text{background})\\
    \;\text{arbitrary warp–bubble profiles}, & r \le R
  \end{cases}
\]
and each profile function is smooth on $[0,R]$, with
\[
  A(0),B(0),C(0),D(0) \;\text{finite}, \quad
  A(R)=B(R)=C(R)=D(R)=1.
\]

\section{Worked Example (Sketch)}
Starting from the general metric above:
\begin{itemize}
  \item \emph{Eliminate cross-terms:}
    $\;g_{t\phi}=g_{r\theta}=0$ by staticity + symmetry.
  \item \emph{Drop $\theta$–dependence:}
    Functions become $g_{\mu\nu}(r)$ only.
  \item \emph{Restrict radial support:}
    $g_{\mu\nu}(r)=\eta_{\mu\nu}$ for $r>R$.
  \item \emph{Rename profiles:}
    $g_{tt}=-A(r)$, $g_{rr}=B(r)$, $g_{\theta\theta}=C(r)\,r^2$, $g_{\phi\phi}=D(r)\,r^2\sin^2\theta$.
\end{itemize}
Thus the simplified ansatz cleanly separates background from warp-bubble interior and makes all subsequent curvature computations depend only on four radial functions.

\end{document}
